% ============================================================================
% ARCHITECTURE
% ============================================================================

\section{Architecture}
\label{sec:architecture}

\subsection{High-Level Overview}

The \ctx{} architecture follows a modular design with clear separation between context generation and code intelligence features. The system is organized into distinct layers, each with well-defined responsibilities.

\begin{figure}[h]
\centering
\begin{tikzpicture}[
    node distance=0.5cm,
    box/.style={rectangle, draw, rounded corners, minimum width=3cm, minimum height=0.8cm, align=center, fill=blue!10},
    bigbox/.style={rectangle, draw, rounded corners, minimum width=6.5cm, minimum height=2.5cm, align=center},
    layer/.style={rectangle, draw, dashed, rounded corners, fill=gray!5},
    arrow/.style={-{Stealth[length=2mm]}, thick}
]

% CLI Layer
\node[box, fill=orange!20, minimum width=14cm] (cli) {CLI Layer (\file{main.rs}, \file{cli.rs})};

% Two main subsystems
\node[bigbox, below=0.8cm of cli, xshift=-3.5cm, fill=green!10] (context) {};
\node[above=0.1cm] at (context.north) {\textbf{Context Generation}};

\node[bigbox, below=0.8cm of cli, xshift=3.5cm, fill=purple!10] (intel) {};
\node[above=0.1cm] at (intel.north) {\textbf{Code Intelligence}};

% Context Generation components
\node[box, fill=green!20] at ([yshift=-0.3cm]context.center) (walker) {\file{walker.rs}};
\node[box, fill=green!20, left=0.2cm of walker, yshift=-0.8cm] (tree) {\file{tree.rs}};
\node[box, fill=green!20, right=0.2cm of walker, yshift=-0.8cm] (formatter) {\file{formatter.rs}};

% Code Intelligence components
\node[box, fill=purple!20] at ([yshift=0.2cm]intel.center) (index) {\file{index/mod.rs}};
\node[box, fill=purple!20, below=0.3cm of index] (parser) {\file{parser/mod.rs}};
\node[box, fill=purple!20, left=0.1cm of parser, yshift=-0.8cm] (db) {\file{db/schema.rs}};
\node[box, fill=purple!20, right=0.1cm of parser, yshift=-0.8cm] (analytics) {\file{analytics/}};

% Storage Layer
\node[box, fill=yellow!20, below=3.5cm of cli, minimum width=5cm] (sqlite) {SQLite (\code{.ctx/codebase.sqlite})};
\node[box, fill=yellow!20, right=0.5cm of sqlite, minimum width=4cm] (duckdb) {DuckDB (In-Memory)};

% Arrows
\draw[arrow] (cli) -- (context);
\draw[arrow] (cli) -- (intel);
\draw[arrow] (index) -- (parser);
\draw[arrow] (db) -- (sqlite);
\draw[arrow] (analytics) -- (duckdb);
\draw[arrow, dashed] (duckdb) -- (sqlite);

\end{tikzpicture}
\caption{High-level architecture diagram}
\label{fig:architecture}
\end{figure}

\subsection{Storage Strategy}

The system employs a dual-database architecture:

\subsubsection{SQLite --- Persistent Storage}

All persistent data resides in a single SQLite file at \code{.ctx/codebase.sqlite}:

\begin{lstlisting}[style=sqlcode, caption={Core database schema}]
-- Files table
CREATE TABLE files (
    path TEXT PRIMARY KEY,
    content_hash TEXT NOT NULL,
    size_bytes INTEGER,
    language TEXT,
    last_indexed INTEGER,
    source_compressed BLOB  -- gzip compressed
);

-- Symbols table
CREATE TABLE symbols (
    id TEXT PRIMARY KEY,        -- "file_path::name::line"
    name TEXT NOT NULL,
    kind TEXT NOT NULL,         -- function, struct, enum, etc.
    file_path TEXT NOT NULL,
    line_start INTEGER,
    line_end INTEGER,
    visibility TEXT DEFAULT 'private',
    signature TEXT,
    brief TEXT,                 -- First line of docstring
    parent_id TEXT,
    FOREIGN KEY (file_path) REFERENCES files(path)
);

-- Edges table (call graph)
CREATE TABLE edges (
    source_id TEXT NOT NULL,    -- Caller symbol ID
    target_name TEXT NOT NULL,  -- Called function name
    target_id TEXT,             -- Resolved symbol ID (if found)
    kind TEXT NOT NULL,         -- 'calls', 'imports', etc.
    line INTEGER,
    context TEXT,               -- Code snippet
    FOREIGN KEY (source_id) REFERENCES symbols(id)
);
\end{lstlisting}

\textbf{Benefits of SQLite:}
\begin{itemize}
    \item Single file for easy management, backup, and deletion
    \item ACID transactions for safe concurrent access
    \item Built-in FTS5 for full-text search
    \item Widely supported with excellent tooling
    \item Optimized for insert-heavy indexing workloads
\end{itemize}

\subsubsection{DuckDB --- Analytical Queries}

For complex analytical queries requiring recursive CTEs and graph traversal, an in-memory DuckDB instance attaches to the SQLite database:

\begin{lstlisting}[style=rustcode, caption={DuckDB integration pattern}]
// Open in-memory DuckDB
let conn = Connection::open_in_memory()?;

// Attach SQLite database read-only
conn.execute(
    "ATTACH 'codebase.sqlite' AS code (TYPE sqlite, READ_ONLY)",
    [],
)?;

// Run recursive graph queries
conn.execute("
    WITH RECURSIVE call_graph AS (
        SELECT source_id, target_id, 1 as depth
        FROM code.edges WHERE source_id = ?
        UNION ALL
        SELECT e.source_id, e.target_id, cg.depth + 1
        FROM code.edges e
        JOIN call_graph cg ON e.source_id = cg.target_id
        WHERE cg.depth < ?
    )
    SELECT * FROM call_graph
", params)?;
\end{lstlisting}

\textbf{Benefits of DuckDB:}
\begin{itemize}
    \item Columnar storage efficient for analytical queries
    \item Native support for recursive CTEs
    \item Direct SQLite attachment without ETL
    \item In-memory operation with no additional files
    \item OLAP-optimized for aggregations and joins
\end{itemize}

\subsection{Parsing Strategy}

The system uses Tree-sitter for parsing, providing:

\begin{itemize}
    \item \textbf{Speed} --- Incremental parsing handles large files efficiently
    \item \textbf{Accuracy} --- Real parser, not regex-based heuristics
    \item \textbf{Consistency} --- Same API across all supported languages
    \item \textbf{Error Tolerance} --- Parses incomplete or invalid code gracefully
    \item \textbf{Battle-tested} --- Used by GitHub, Neovim, Zed, and other major tools
\end{itemize}

\subsection{Symbol Identification}

Symbols are identified by a composite key ensuring uniqueness:

\begin{verbatim}
    file_path::name::line
\end{verbatim}

For example:
\begin{verbatim}
    src/auth/handler.ts::handleAuth::45
    src/parser/rust.rs::extract_symbols::200
\end{verbatim}

This scheme handles:
\begin{itemize}
    \item Same function name in different files
    \item Overloaded functions
    \item Nested functions and closures
\end{itemize}

\subsection{Incremental Indexing}

File changes are tracked using content hashes rather than timestamps:

\begin{lstlisting}[style=rustcode, caption={Incremental update detection}]
fn needs_update(&self, path: &str, content: &str) -> bool {
    let new_hash = sha256(content);
    let old_hash = self.db.get_hash(path);
    new_hash != old_hash
}
\end{lstlisting}

This approach ensures:
\begin{itemize}
    \item Only changed files are re-parsed
    \item Unchanged files are skipped instantly
    \item More reliable than timestamp-based detection
    \item Handles file moves and renames correctly
\end{itemize}
